\label{sec:overview}

This section describes the design and implementation of \system. The goal of \system is to add a persistent mistake-memory layer to an LLM vulnerability detection pipeline. A standard baseline model analyzes each code sample independently. In contrast, \system stores previous incorrect predictions in a database and retrieves similar mistakes when analyzing a new sample. This lets the model receive explicit warnings about similar cases it has failed on before.

\subsection{System Architecture}

\system consists of four main components: the baseline LLM predictor, the mistake extraction module, the SQLite mistake database, and the retrieval-augmented prediction module. The baseline predictor receives a C/C++ code sample and asks the LLM to return a vulnerability prediction in the benchmark's required JSON format. The prediction is then compared against the ground truth label. If the prediction is incorrect, the mistake extraction module converts the failure into a structured record and inserts it into the mistake database. During future inference, the retrieval module searches the database for similar past mistakes and adds the retrieved records to the prompt before asking the same LLM to make a new prediction.

The overall workflow is:

\begin{enumerate}
    \item Run the base LLM on labeled benchmark samples.
    \item Parse the model output into predicted vulnerability label, CWE, and vulnerable lines.
    \item Compare the prediction with the ground truth.
    \item Store incorrect predictions as structured mistake records in SQLite.
    \item For a new input sample, retrieve similar stored mistakes using TF-IDF similarity.
    \item Insert the retrieved mistakes into the prompt as cautionary examples.
    \item Run the LLM again using the ErrorLock-augmented prompt.
\end{enumerate}

This architecture separates the base model from the mistake-memory layer. \system does not fine-tune the LLM or change its internal parameters. Instead, it changes the surrounding inference pipeline by adding database-backed memory and retrieval-augmented prompting.

\subsection{Baseline Vulnerability Detection}

In the original CASTLE and Juliet experiments, the baseline component uses Qwen2.5-Coder-7B-Instruct as the base LLM. Qwen2.5 includes specialized models for coding under the Qwen2.5-Coder family, including 1.5B and 7B model sizes~\cite{qwen25coder}. In our experiments, each input prompt contains the benchmark instruction and the source code sample. The model is asked to output a JSON list describing detected vulnerabilities. An empty list means the model predicts that the code is safe. A non-empty list means the model predicts that the code contains at least one vulnerability.

After generation, the parser extracts three pieces of information from the model output: the binary vulnerability label, the predicted CWE, and the predicted vulnerable lines. If the output cannot be parsed as the required JSON format, the prediction is marked as invalid. Parsed predictions are then compared with the benchmark ground truth. This baseline stage serves two purposes. First, it gives a normal LLM-only result for comparison. Second, it produces the mistakes that are inserted into the ErrorLock database.

\subsection{Mistake Record Schema}

Each incorrect prediction is converted into a structured mistake record. The database stores both ground-truth information and model-output information. The main fields are:

\begin{itemize}
    \item \textbf{sample\_id}: the identifier of the benchmark sample.
    \item \textbf{source\_name}: the source split or benchmark used to build the database.
    \item \textbf{code}: the original source code sample.
    \item \textbf{true\_label}: the ground-truth binary vulnerability label.
    \item \textbf{pred\_label}: the model's predicted binary vulnerability label.
    \item \textbf{true\_cwe}: the ground-truth CWE category.
    \item \textbf{pred\_cwe}: the model's predicted CWE category, if available.
    \item \textbf{true\_lines}: the ground-truth vulnerable lines, when provided by the benchmark.
    \item \textbf{pred\_lines}: the model's predicted vulnerable lines.
    \item \textbf{error\_type}: a label describing the mistake, such as false positive or false negative.
    \item \textbf{raw\_output}: the original raw text generated by the model.
\end{itemize}

The error type is computed from the true and predicted labels. If the true label is safe but the model predicts vulnerable, the mistake is stored as a false positive. If the true label is vulnerable but the model predicts safe, the mistake is stored as a false negative. Other parsed incorrect outputs are stored as wrong predictions. This schema makes the mistake database inspectable because each record preserves both what the model predicted and what the correct answer should have been.

\subsection{Database Backend}

The ErrorLock prototype uses SQLite as the backend database engine. SQLite is a good fit for this project because it is self-contained, serverless, zero-configuration, and stores the database locally in a file~\cite{sqlite}. This makes it easy to use inside a Jupyter or Colab notebook without setting up a separate database server. Each model receives its own SQLite mistake database, so the stored mistakes are model-specific. This is important because different LLMs may make different errors on the same benchmark.

For a larger production system, SQLite could be replaced by a client-server database such as PostgreSQL or by a vector database that supports embedding-based similarity search. However, for this course-project prototype, SQLite is sufficient because the experiments use small benchmark subsets and the main goal is to test whether persistent mistake memory can influence future predictions.

\subsection{Mistake Retrieval}

After building the mistake database, \system retrieves similar mistakes for each new input sample. The current prototype uses TF-IDF similarity over a textual representation of each mistake record. TF-IDF is a standard text representation that weights terms according to how often they appear in a document and how rare they are across the corpus. We use scikit-learn's implementation to vectorize mistake records and query samples~\cite{sklearn}.

For each stored mistake, \system builds retrieval text containing the sample identifier, error type, true label, predicted label, true CWE, predicted CWE, and source code. For a new test sample, the query text contains the sample code and its CWE category. The retriever computes cosine similarity between the query vector and each mistake vector, then selects the top-$k$ most similar mistakes. In our experiments, we use $k=3$. Cosine similarity measures similarity as the normalized dot product between two vectors, which makes it suitable for comparing TF-IDF vectors~\cite{sklearn}.

This retrieval method is simple and lightweight. It does not require training a separate embedding model or indexing service. However, it also has limitations. TF-IDF mainly captures lexical similarity, so it may miss semantically similar code patterns written with different variable names or structure. We discuss this limitation further in Section~\ref{sec:conclusion}.

For the supplemental OpenRouter experiment, we also implement a compact RAG memory variant. Instead of sending the entire mistake table to the model, the system first retrieves a larger candidate pool from SQLite, reranks candidates using similarity and prediction-pattern signals, diversifies the selected evidence by CWE and error type, and sends only a small evidence pack to the prompt. In the reported setting, the RAG retriever considers 25 candidate mistakes and selects 6 evidence rows. This keeps the prompt much smaller than all-memory retrieval while still giving the model broader context than the original top-$k$ retrieval.

\subsection{ErrorLock-Augmented Prompting}

The retrieved mistakes are inserted into the prompt before the new code sample. Each retrieved record is formatted as a past mistake with its similarity score, error type, ground-truth label, ground-truth CWE, predicted label, predicted CWE, and a short lesson. For example, if the retrieved mistake was a false positive, the prompt warns the model not to over-report similar safe code patterns. If the retrieved mistake was a false negative, the prompt warns the model not to ignore similar risky patterns.

The prompt also tells the model not to blindly copy the retrieved labels. This is important because the retrieved mistakes are only guidance. The current code sample may be similar to a past mistake but still have a different correct label. Therefore, \system uses retrieved mistakes as cautionary examples rather than as direct answers.

\subsection{Trace and Inspectability}

The benchmarks used in this project provide structured labels, CWE categories, and in some cases vulnerable-line information, but they do not provide full reasoning traces explaining why each sample is vulnerable or safe. \system therefore creates its own lightweight trace during inference. For each ErrorLock prediction, the output CSV stores the predicted label, predicted CWE, retrieved mistake identifiers, retrieved error types, retrieved true CWEs, retrieved predicted CWEs, and retrieval similarity scores.

This trace is useful for analysis because it lets us inspect which past mistakes influenced each prediction. For example, if ErrorLock reduces false positives but increases false negatives, we can examine whether the retrieved records were mostly false-positive mistakes that encouraged the model to become more conservative. This makes the system more transparent than a plain LLM baseline, where each prediction is produced without an explicit record of relevant past failures.

\subsection{Implementation Details}

We implemented \system in Python in a Jupyter/Google Colab notebook. The implementation uses Hugging Face Transformers to load Qwen2.5-Coder-7B-Instruct, SQLite for mistake storage, pandas for data processing, and scikit-learn for TF-IDF vectorization and cosine-similarity retrieval. The artifact repository contains the notebook, generated CSV result files, and SQLite mistake databases.

The implementation is organized around reusable functions for model loading, baseline prediction, output parsing, mistake database construction, retrieval, ErrorLock prompt construction, ErrorLock prediction, and metric computation. The original experiment pipeline runs the baseline and ErrorLock variants on CASTLE and Juliet, then exports report-ready summary tables. The supplemental OpenRouter RAG run uses the same database idea, but evaluates a stronger API model with the compact evidence-pack retriever on the full CASTLE-C250 split.
